%-----------------------------------------------------------------------------
% Posłowia
%-----------------------------------------------------------------------------

\cleardoublepage
\chapter*{Posłowie}

\phantomsection

\lettrine{Z}{nana} i~dobrze udokumentowana historia rodziny Gomulskich z~Desna
rozpoczyna się od~narodzin Jana Gomoły w~1789 roku w~Mysłowie, na~ziemi
łukowskiej ówczesnego województwa lubelskiego. Od~tego momentu przez kolejne
sto lat Jan Gomoła oraz jego potomkowie szukali swojego miejsca na~ziemi,
przeprowadzając się coraz dalej na~północny zachód od~rodzinnych stron~--
zbliżając się nieuchronnie ku~Warszawie. Mieszkali kolejno w~Rudzienku,
Zamieniu, Grzebowilku, Anielinie, Wólce Mińskiej, Mińsku Mazowieckim,
Mikanowie oraz~Karolinie, aż~w~końcu pod~koniec XIX~wieku znaczna część
potomków Jana Gomoły osiadła na~stałe na~terenie obecnego Desna, a~ich
rodziny żyją tam już od~przeszło stu trzydziestu lat.

Wraz z~tą rodzinną wędrówką ewoluowało również rodowe nazwisko. Według
przedstawionej w~książce hipotezy~-- od~przydomków szlacheckich rodziny
Dziewulskich~-- drobnej szlachty zagrodowej północnej Lubelszczyzny
(\emph{Gomolik}, \emph{Gomolak}, \emph{Gomoła}, \emph{Gomuła})~-- przez
pierwsze, półoficjalne protonazwiska (\emph{Gomolak}, \emph{Gomoła}
oraz~\emph{Gumulski}), aż~do~obecnego, pełnoprawnego brzmienia nazwiska
Gomulski, które utrwaliło się na~dobre w~dokumentach urzędowych dopiero
po~odzyskaniu przez Polskę niepodległości, a~dokładnie pod~koniec lat
dwudziestych XX~wieku.

Zmieniał się również główny zawód wykonywany przez członków rodziny
Gomulskich. Jan Gomoła, jego synowie Piotr Franciszek i~Andrzej, a~także syn
Andrzeja, Feliks, byli wiejskimi kowalami. Pełnili tę rolę przez około sto lat
kolejno w~Rudzienku, Zamieniu, Grzebowilku oraz~Anielinie, aż~do~lat
dwudziestych XX~wieku. Dalsi potomkowie Jana Gomoły utrzymywali się już
przede wszystkim z~uprawy ziemi. Taka sytuacja miała miejsce aż~do~połowy
XX~wieku, kiedy to~pracę na~roli zaczęli łączyć z~zatrudnieniem w~zakładach
przemysłowych powojennej Warszawy. Praca poza gospodarstwem stopniowo zaczęła
dominować nad obowiązkami w~gospodarstwie, aż~w~końcu, na~początku XXI~wieku,
już tylko nieliczni potomkowie Jana Gomoły trudnili się rolnictwem w~celach
innych niż~wyłącznie hobbystyczne.

Przez ostatnie trzysta lat istotnej ewolucji uległ również status społeczny
członków rodziny Gomulskich~-- od~przypuszczalnych korzeni wśród zubożałej
drobnej szlachty zagrodowej ziemi łukowskiej, poprzez chłopskich rzemieślników
wschodniej części Mazowsza, wolnych chłopów środkowego Mazowsza
i~podwarszawskich chłoporobotników, aż~po~pełnoprawnych przedstawicieli
polskiej klasy średniej zamieszkujących wschodnie przedmieścia stolicy.
Historia rodziny Gomulskich, jak~w~soczewce, pokazuje, jak bardzo zmieniła się
Polska na~przestrzeni ostatnich trzech wieków, przy~czym najsilniejszymi
katalizatorami tych zmian wydają~się być przede wszystkim dwa procesy
historyczne: uwłaszczenie ludności chłopskiej w~drugiej połowie XIX~wieku
oraz~stopniowe odzyskiwanie przez Polskę niepodległości na~początku XX~wieku.
W~tym okresie wiele rodzin chłopskich, w~tym również rodzina Gomulskich,
zyskało własną ziemię i~stałe miejsce zamieszkania oraz dostęp do~pierwszych
zrębów edukacji dla swoich dzieci. We~wcześniejszych pokoleniach większość
przodków oraz potomków Jana Gomoły była analfabetami, niepotrafiącymi się
najczęściej podpisać nawet pod~aktem własnego ślubu. Suma tych dwóch procesów
historycznych przyczyniła się do~poprawy perspektyw zawodowych kolejnych
pokoleń, wzrostu poziomu ich wykształcenia oraz awansu społecznego, które
trwają, na~szczęście, nieprzerwanie aż~do~dziś.

\emph{Desna} na przełomie XIX i~XX~wieku stała się więc stabilnym domem
dla~dużej części potomków Jana Gomoły~-- pierwszym takim miejscem od~wielu
pokoleń, po~tym, jak~-- według przedstawionej w~książce hipotezy~-- jego
przodkowie opuścili rodzinne Dziewule. Domem, który nie~był im~przez nikogo
dany, ale~który sami wybudowali i~wyposażyli dzięki ciężkiej pracy oraz
wielkiemu poświęceniu, miesiąc po~miesiącu spłacając zobowiązania zaciągnięte
wspólnie z~nowymi sąsiadami w~ramach spółki \emph{Desna}. Zobowiązania, które
pozwoliły im~nie tylko nabyć własne miejsce na~ziemi, ale również~-- a~być
może przede wszystkim~-- w~większym stopniu decydować o~własnym losie.

Gdy zaczynałem pisać tę książkę, chciałem zaspokoić dziecięcą ciekawość,
usystematyzować wiedzę o~rodzinie Gomulskich z~Desna, a~także w~przystępny
sposób przekazać ją~kolejnym pokoleniom. Nie~spodziewałem się, że~odkryję tak
wiele fascynujących historii, które nie~tylko pozwoliły mi~lepiej zrozumieć
przodków, ale także dostrzec związek ich codziennych wyborów z~historią mojej
ojczyzny. Mam nadzieję, że~książka ta~będzie dla~mojej rodziny cennym źródłem
wiedzy o~naszych korzeniach oraz~inspiracją do~dalszego poznawania naszych
dziejów. Wierzę również, że~okaże się~ona interesująca dla~czytelników spoza
rodziny Gomulskich~-- że~udało mi się pokazać, iż~nie~trzeba mieć drzewa
genealogicznego wypełnionego sławnymi postaciami, by~móc opowiedzieć ciekawą
historię o~jednej z~wielu zwyczajnych polskich rodzin oraz~czuć dumę
i~wdzięczność, że~jest się częścią tej opowieści. \\

\begin{flushright}
\emph{Mateusz Gomulski, 10~października 2026~r.} \\
\emph{Desno~9a, Drugi Dom Rodzinny Autora} \\
\emph{Po Pozytywnym Zakończeniu Leczenia Onkologicznego} 
\end{flushright}